\documentclass[11pt,a4paper]{article}
\usepackage[utf8]{inputenc}
\usepackage[T1]{fontenc}
\usepackage{amsmath,amssymb,amsthm}
\usepackage{geometry}
\usepackage{booktabs}
\usepackage{enumitem}
\usepackage{hyperref}
\usepackage{verbatim}

\geometry{margin=1in}

\newtheorem{theorem}{Theorem}
\newtheorem{lemma}[theorem]{Lemma}
\newtheorem{proposition}[theorem]{Proposition}
\newtheorem{definition}{Definition}
\newtheorem{corollary}[theorem]{Corollary}
\newtheorem{conjecture}{Conjecture}
\newtheorem{question}{Open Question}

\title{Breaking Bellman:\\[6pt]
\large Null Uncertainty as a Computable Structural Boundary\\Where Value Functions Fail}
\author{Ron Itelman\\Chair, W3C Context Graph Community Group\\\texttt{ron@ronitelman.com}}
\date{Working Paper --- April 2026\\[4pt]
\small Not for citation. Submitted for open review and falsification.}

\begin{document}
\maketitle

\begin{abstract}
The Bellman equation is the foundation of sequential decision-making in AI, 
economics, and control theory. It has one precondition that, to my knowledge, has never been explicitly formalized: that the decision-maker can observe the full state 
relevant to its decision. 

This paper is not about the equation producing the 
wrong answer. It is about identifying the structural conditions under which 
the state space the equation operates over becomes fictitious --- rendering 
the equation's operation meaningless before a single value is computed.

We identify four failure modes at the serialization boundary between agents, 
each with distinct epistemic consequences. We introduce a four-category scoring framework 
for Layer~1 agent behavior: \textbf{Asked} ($A$): the agent requested 
clarification before proceeding; \textbf{Detected} ($D$): the agent named the 
ambiguity but assumed a default and proceeded; \textbf{Undetected} ($U$): the 
agent silently assumed with no record of ambiguity; \textbf{Hallucinated} ($H$): 
the agent introduced a dimension absent from both the input and the world state. 
A critical non-obvious result: $D$ and $U$ are behaviorally distinct at Layer~1 
but produce identical epistemic consequences at Layer~2 --- both result in null 
uncertainty, the absence of a slot for the dimension in the downstream state 
space. Only $A$ prevents null uncertainty. $H$ is strictly worse than either: 
it expands the downstream state space with fictitious dimensions over which the 
Bellman equation optimizes, and no rotation can repair it because no ground 
truth exists anywhere in the pipeline.

The mechanism is serialization: when one agent writes its assumptions into 
JSON, ambiguity collapses into false certainty, provenance is severed, and 
the downstream agent's state space is simultaneously truncated of real 
dimensions and expanded with fabricated ones. We provide a minimal empirical 
protocol with a hardcoded answer key computable before any agent runs: 7.22 
bits of boundary entropy, 4 required cross-boundary measurements, 150 
indistinguishable configurations. We show the failure is not fixable by 
improving the agent, because the information needed to detect null dimensions 
is on the other side of the boundary, and hallucinated dimensions have no 
other side. We position this result against POMDPs, HMS unawareness, model 
misspecification, and Dec-POMDPs, and show that none addresses the specific 
failure produced by serialization. We pose six open questions, including 
whether a hallucination-adjusted coverage metric $C^*$ is computable from 
outside the pipeline.
\end{abstract}

\section{Failure Modes in Systems of Decision Networks}
\subsection{Where Bellman Breaks: Scope and Motivation}

The Bellman equation is embedded in every reinforcement learning 
algorithm, every dynamic programming solution, and every AI agent 
that plans over multiple steps. Its precondition --- that the agent 
can observe the state relevant to its decision --- is so fundamental 
it is rarely stated. This paper is about what happens when that 
precondition is not merely violated by noise or partial observation, 
but structurally destroyed by the mechanism through which agents 
communicate.

The accelerated deployment of AI agents in enterprise workflows has 
created a new failure surface. When agents operate in pipelines --- 
one agent's output becoming the next agent's input --- each 
serialization boundary is a point at which the precondition can fail 
silently. The downstream agent receives well-formed data, passes all 
syntactic validation, and optimizes confidently over a state space 
that may be simultaneously missing real dimensions and populated with 
fictitious ones. No internal diagnostic detects this. The value 
function converges --- to the wrong number.

This failure is not a property of any particular agent, model, or 
prompt. It is a property of serialization itself: the act of writing 
assumptions into structured data collapses ambiguity into false 
certainty, severs provenance, and truncates the downstream agent's 
state space. The failure compounds with every agent hop, every nested 
assumption, and every repetition over time.

The W3C Context Graph Community Group is developing formal 
infrastructure to make this failure mode measurable and engineerable, 
in the same way Shannon's information theory made communication 
capacity measurable and engineerable~\cite{shannon1948}. This paper 
establishes the mathematical boundary conditions: where exactly 
Bellman breaks, why it cannot be fixed by improving the agent, and 
what the minimum viable architectural response looks like. The 
mathematics in this working paper is submitted for open review and 
independent verification.

\subsection{Risk Surface Management at System Boundaries}

The distinction between hallucination and null uncertainty is not merely 
epistemic --- it determines whether the failure mode is engineerably 
controllable.

\textbf{Hallucination is enumerable in advance.} Within a system, a 
designer can bound the set of dimensions an agent might fabricate, because 
the agent's generative behavior is a function of its training and 
architecture --- both of which are inspectable. A compliance envelope can 
be defined:

\begin{equation}
    \mathcal{K}_{\text{permitted}} = \{k_1, k_2, \ldots, k_m\}
\end{equation}

Any $k \notin \mathcal{K}_{\text{permitted}}$ appearing in the agent's 
output is detectable by schema validation before it crosses a boundary. 
This is standard deterministic systems engineering: define the output 
contract, reject violations at the gate. Hallucination risk, while real, 
is therefore \emph{bounded and controllable within a system}.

\textbf{Null uncertainty is not enumerable in advance.} The defining 
property of a null variable is that it has no slot in the agent's state 
space --- which means the agent cannot report it, the designer cannot 
anticipate it from within the system, and no schema validation can flag 
its absence. For the field \texttt{account: TRD-4471}, the system can 
enumerate permitted output dimensions. It cannot enumerate the dimensions 
of the codebook the sender used to produce that value, because those 
dimensions exist on the other side of the boundary:

\begin{equation}
    |\Omega_{\text{null}}| \text{ is not computable from within } S
\end{equation}

A permitted output schema for \texttt{account} might read:
\begin{verbatim}
"account": ["opaque_identifier"]
\end{verbatim}
This prevents hallucination of \texttt{account\_authorization}. It does 
nothing to surface whether \texttt{TRD-4471} refers to an equity account, 
a derivatives account, a custodial account, or a routing identifier --- 
because those distinctions live in the sender's codebook, not in the 
received value.

The asymmetry is structural:
\begin{center}
\begin{tabular}{lcc}
\toprule
\textbf{Failure Mode} & \textbf{Enumerable in advance} & 
\textbf{Controllable within system} \\
\midrule
Hallucination ($H$)  & Yes & Yes --- schema validation \\
Null uncertainty     & No  & No --- requires cross-boundary rotation \\
\bottomrule
\end{tabular}
\end{center}

This asymmetry is why null uncertainty is the harder problem, and why 
the engineering response must be architectural rather than 
parametric. No improvement to the agent resolves null uncertainty, 
because the agent would need to improve by knowing what it does not 
know --- which is precisely what null uncertainty forecloses.

\textbf{The mechanics are general.} Boundaries of this kind exist not 
only between organizations but within them: between teams, between 
microservices, between the agent that executes a trade and the agent 
that assesses the resulting risk. The failure mechanism --- 
serialization collapsing ambiguity into false certainty, severing 
provenance, truncating the downstream state space --- is identical 
regardless of where the boundary falls. The entropy bound 
(Section~\ref{sec:answer}) is the same calculation whether the 
boundary is internal or external. The codebook problem is the same 
problem.

\textbf{Rotation has a cost.} Each cross-boundary measurement --- 
each act of shifting orientation from inward (own codebook) to 
outward (cross-boundary signal) and recording the return --- requires 
time, compute, or attention. This cost is non-eliminable: there is no 
way to verify codebook alignment without acquiring information from 
the other side of the boundary. The question is not whether the cost 
exists but who or what pays it.

\begin{center}
\begin{tabular}{lccc}
\toprule
\textbf{Boundary type} & \textbf{Math} & \textbf{Rotation possible} & 
\textbf{Who pays the cost} \\
\midrule
Internal, low volume  & Identical & Yes & Human, manually \\
Internal, high volume & Identical & Yes & Requires instrumentation \\
External, any volume  & Identical & Yes & Requires instrumentation \\
\bottomrule
\end{tabular}
\end{center}

At low volume and within a single organization, a skilled engineer 
can perform rotations manually --- asking what a column means, 
checking the timezone, verifying the currency convention. This is 
what ETL engineers do. It is informal, unaudited, and does not scale. 
At agent-paced operation, where thousands of boundary crossings occur 
per hour without human intervention, manual rotation is not a 
viable architecture.

Across external system boundaries, manual rotation is additionally 
constrained by jurisdiction: the engineer on one side does not have 
authority to inspect the codebook on the other side, and may not 
know the boundary exists. The information needed is structurally 
inaccessible without a protocol both sides have agreed to participate 
in.

This is the engineering motivation for the \emph{Observatron}: a 
stateful boundary node that rotates between inward and outward 
orientations at a system boundary, records the return as a canonical 
claim, and converts null uncertainty into defined uncertainty before 
serialization occurs~\cite{itelman2026lh}. The Observatron does not 
change the mathematics. It makes the cost of rotation payable at 
scale, in a form that is auditable, machine-readable, and composable 
across pipelines. The boundary math is the same. What changes is 
whether anyone is paying the rotation cost at all.
%% ============================================================
\section{Background: The Bellman Equation}
%% ============================================================

In 1957, Richard Bellman published \emph{Dynamic Programming}, establishing the mathematical foundation for sequential decision-making under uncertainty. The core insight is that any multi-step decision problem can be decomposed into a chain of one-step decisions. The Bellman equation expresses this recursively:
\[
V^*(s) = \max_{a \in A} \left[ R(s, a) + \gamma \sum_{s'} P(s' | s, a)\, V^*(s') \right]
\]
where $V^*(s)$ is the value of being in state $s$ under the optimal policy, $R(s, a)$ is the immediate reward for taking action $a$ in state $s$, $\gamma$ is a discount factor, and $P(s'|s,a)$ is the probability of transitioning to state $s'$.

The equation says: the value of your current situation is the best action you can take now, plus the discounted value of where that action puts you next. This is the basis of every reinforcement learning algorithm, every dynamic programming solution, every AI agent that plans over multiple steps.

The equation has one precondition that appears so obvious it is rarely stated: \emph{the agent must be able to observe $s$}. If the agent cannot observe the state, it cannot evaluate which action is best. The equation produces a number, but that number does not correspond to the true value of the situation.

This paper is about what happens when that precondition fails --- specifically, at the boundaries between systems where data crosses but codebooks do not.


%% ============================================================
\section{The Precondition}
%% ============================================================

\begin{definition}[True State]
At a system boundary, the true state is $s_{\mathrm{true}} = (s, \omega)$, where $s$ is the data the agent receives and $\omega \in \Omega$ is the codebook configuration --- the date format, unit convention, semantic definition, and reference frame under which $s$ should be interpreted.
\end{definition}

\begin{definition}[Observation Function]
The observation function at a system boundary is $O(s, \omega) = s$. The data crosses the boundary. The codebook does not. The agent sees $s$ but not $\omega$.
\end{definition}

\begin{definition}[Resolution Capacity]
The resolution capacity $R$ is the number of configurations in $\Omega$ that the agent can distinguish from $s$ alone. For a bare value with no metadata, $R = 1$: the agent can read the value but cannot determine which codebook produced it.
\end{definition}

\noindent When $|\Omega| > R$, the pigeonhole principle guarantees that at least $|\Omega| - R$ configurations produce identical observations. The agent cannot distinguish them. The Bellman equation treats them as the same state. They are not the same state.

\section{The Observability Condition}

We state the core observability result as a lemma rather than a 
theorem. It is cleared ground --- a direct consequence of the 
projection structure of the observation function --- not the primary 
contribution of this paper. The primary contributions are: (1) the 
identification of serialization as the physical mechanism that 
produces this condition at system boundaries, (2) the four-category 
scoring framework (A, D, U, H) that characterizes how agents fail 
under this condition, and (3) the H category and its consequence that 
rotation-based recovery is undefined for hallucinated dimensions. The 
lemma establishes why recovery is impossible within a fixed state 
space; the rest of the paper establishes what that means in practice 
and at scale.

\begin{lemma}[Bellman Observability Failure]
\label{thm:bellman}
If:
\begin{enumerate}[label=(\roman*)]
  \item there exist $\omega_1, \omega_2 \in \Omega$ such that the 
  optimal action differs:
  $\arg\max_a R(s, \omega_1, a) \neq \arg\max_a R(s, \omega_2, a)$,
  \item $O(s, \omega_1) = O(s, \omega_2) = s$,
\end{enumerate}
then the observed value function $\hat{V}(s) \neq V^*(s, \omega)$ 
for at least one $\omega$, and no function computable from $s$ alone 
can restore equality.
\end{lemma}

\begin{proof}
Different optimal actions yield different optimal values: 
$V^*(s, \omega_1) \neq V^*(s, \omega_2)$. But $\hat{V}(s)$ returns 
one number. One number cannot equal two different values. Every 
function of $s$ returns the same output for $\omega_1$ and $\omega_2$ 
because $O$ projects away $\omega$. Therefore no refinement within 
the agent's state space recovers the true value function. \qed
\end{proof}

The divergence gap --- the minimum guaranteed loss --- is:
\[
\Delta = |V^*(s, \omega_1) - V^*(s, \omega_2)| > 0
\]

\noindent This result is closely related to the unawareness 
literature, particularly Heifetz, Meier, and Schipper 
(HMS)~\cite{hms2006}, who formalize that agents inhabiting a state 
space $S_\alpha$ cannot express propositions requiring a richer space 
$S_\beta \succ S_\alpha$. The present result is an application of 
that structure to serialization pipelines. What we add beyond HMS is 
threefold: the physical mechanism (serialization) that produces the 
lattice, a computable bound on the exposure (Section~\ref{sec:answer}), 
and the H category (Section~\ref{sec:hallucination}), which has no 
analogue in HMS --- HMS addresses missing propositions, not fabricated 
dimensions with undefined rotation functions.




%% ============================================================
\section{Where Exactly Does It Break?}
\label{sec:where}
%% ============================================================

The key question is: at which point in a multi-agent workflow does the Bellman equation become invalid? The answer is precise, and it is not where most people would expect.

\subsection{Layer 1: The Trade Agent}

An agent receives this CSV:
\begin{verbatim}
date,amount,location,instrument,direction,order_type,account
4/5/2026,1000000,Singapore,SGX,buy,market,TRD-4471
\end{verbatim}


At this point, the Bellman equation is \emph{already invalid} in the 
sense of Lemma~\ref{thm:bellman}: the agent's state $s$ does not include $\omega$, so its value function is wrong for at least one configuration. However, the agent is at least \emph{proximate} to the boundary. A well-designed agent could, in principle, recognize that the date format is ambiguous and ask. The information it needs exists one hop away --- at the sender.

\textbf{Bellman status at Layer 1:} Invalid, but potentially recoverable.

\subsection{Layer 2: The Risk Agent --- Where It Truly Breaks}

The trade agent outputs structured JSON where every field is concrete: \texttt{"date": "2026-04-05"}, \texttt{"currency": "SGD"}, \texttt{"instrument": "S68.SI"}. A second agent receives this and is asked to assess portfolio risk.

\textbf{This is where Bellman breaks irreparably.} Four things happen simultaneously:

\begin{enumerate}
  \item \textbf{The ambiguity is gone.} Every assumption has been hardened into data. The risk agent has no signal these were ever uncertain.

  \item \textbf{The boundary is no longer adjacent.} The risk agent is two hops from the original sender. It does not know who the sender was, what format they used, or that a format choice was made.

  \item \textbf{The risk agent adds its own assumptions.} Correlation model, volatility window, risk threshold, hedging instrument --- each is a new open boundary with its own $\Omega$.

  \item \textbf{The risk agent's state space is structurally incomplete.} Its state $s_2$ includes Layer 1's output as fact. $\omega_1 \notin s_2$. No function in $\mathcal{F}(S_2)$ can surface the need for it.
\end{enumerate}

\textbf{Bellman status at Layer 2:} Invalid and non-recoverable.

\subsection{The Precise Mechanism}

The break point is the \textbf{serialization boundary} between Layer 1 and Layer 2. When Layer 1 writes its assumptions into JSON:

\begin{enumerate}
  \item \textbf{Entropy collapse:} The date field goes from $H = 1$ bit to $H = 0$ bits. But the collapse was caused by assumption, not measurement. The entropy became invisible.

  \item \textbf{Provenance loss:} The JSON does not record that \texttt{2026-04-05} was derived from an ambiguous input.

  \item \textbf{State space truncation:} $\omega_1$ is not a variable in $S_2$. It is not hidden. It does not exist.
\end{enumerate}


\section{The Geometry of Compounding Null Uncertainty}
\label{sec:compounding}

The previous section identified the serialization boundary as the point where Bellman breaks. This section shows that the failure is not isolated --- it compounds across three independent axes, producing a volume of undetectable exposure that grows with every agent hop, every nested assumption, and every repetition over time.

\subsection{Three Axes}

We define a coordinate system for null uncertainty:

\begin{definition}[Sequential Axis ($X$)]
The sequential axis counts \emph{agent hops}. Each hop $x \in \{1, 2, \ldots, L\}$ is a serialization boundary where one agent's output becomes the next agent's input. At each hop, provenance is severed and the downstream agent's state space is truncated. $L$ is the total depth of the agent pipeline.
\end{definition}

\begin{definition}[Nested Axis ($Y$)]
The nested axis counts \emph{assumptions within a single layer}. At layer $x$, the agent makes $n_x$ assumptions, each with its own interpretation space $\Omega_{x,y}$. The assumptions are stacked: some are inherited from previous layers (invisible, $\omega \notin S_x$) and some are the agent's own (visible within $S_x$). Let $n_x^{\text{inh}}$ be the count of inherited assumptions and $n_x^{\text{own}}$ the count of the agent's own. The total depth at layer $x$ is $n_x = n_x^{\text{inh}} + n_x^{\text{own}}$.
\end{definition}

\begin{definition}[Temporal Axis ($Z$)]
The temporal axis counts \emph{repetitions over time}. Each repetition $z \in \{1, 2, \ldots, T\}$ is an independent execution of the pipeline (e.g., one trade per day over a quarter). At the end of the temporal window, a downstream agent (e.g., a strategy recommender) aggregates the outputs of all $T$ runs.
\end{definition}

\subsection{Entropy at a Point}

At coordinates $(x, y, z)$ --- layer $x$, assumption $y$, run $z$ --- the boundary entropy is:
\[
h(x, y, z) = \begin{cases}
\log_2 |\Omega_{x,y}| & \text{if assumption } y \text{ is Detected ($D$) or Undetected ($U$)} \\
0 & \text{if assumption } y \text{ was Asked ($A$)} \\
\text{excluded} & \text{if assumption } y \text{ is Hallucinated ($H$) --- see Section~\ref{sec:hallucination}}
\end{cases}
\]
\subsection{Entropy per Layer}

At layer $x$ in run $z$, the total entropy is the sum over all assumptions at that layer:
\[
H_x(z) = \underbrace{\sum_{y=1}^{n_x^{\text{inh}}} h(x, y, z)}_{\text{inherited (invisible)}} + \underbrace{\sum_{y=n_x^{\text{inh}}+1}^{n_x} h(x, y, z)}_{\text{own (visible)}}
\]

The critical structural property: the inherited term is \emph{invisible to the agent at layer $x$}. No function in $\mathcal{F}(S_x)$ can compute it. The agent's self-reported uncertainty is only the second sum. The first sum is structurally absent from its state space.

\subsection{Inherited Entropy Grows Monotonically}

At each serialization boundary, the inherited entropy can only increase:
\[
H_{x+1}^{\text{inh}} = H_x^{\text{inh}} + H_x^{\text{own}} = \sum_{k=1}^{x} H_k^{\text{own}}
\]

This follows from the serialization mechanism: every unresolved assumption at layer $x$ is hardened into data at layer $x+1$. Nothing is lost from the inherited set. Everything unresolved in the own set is added to it.

At layer $L$ (the final agent in the pipeline), the total inherited entropy is:
\[
H_L^{\text{inh}} = \sum_{k=1}^{L-1} H_k^{\text{own}}
\]

\textbf{The agent at layer $L$ can see only $H_L^{\text{own}}$. It is blind to $H_L^{\text{inh}}$.}

\subsection{Configuration Space Grows Multiplicatively}

While entropy sums additively, the configuration space grows as a product:
\[
|\Omega_{\text{total}}^{(x)}| = \prod_{k=1}^{x} \prod_{y=1}^{n_k} |\Omega_{k,y}|
\]

For our experiment: Layer 1 has $|\Omega^{(1)}| = 150$. If Layer 2 adds 3 own assumptions with $|\Omega| = 3, 4, 5$ respectively, then:
\[
|\Omega^{(2)}_{\text{total}}| = 150 \times 3 \times 4 \times 5 = 9{,}000
\]

The Layer 2 agent can distinguish at most $\prod R_{2,y} = 1$ configuration from its own state space. At least 8,999 configurations are observationally indistinguishable from verified alignment.

\subsection{Temporal Aggregation}

A strategy agent at the end of a quarter aggregates $T$ pipeline runs. Each run $z$ carries total entropy $H_L(z)$. The strategy agent's total null exposure is:
\[
H_{\text{strategy}} = \sum_{z=1}^{T} H_L(z) = T \cdot \left( \sum_{k=1}^{L} H_k^{\text{own}} \right)
\]

assuming each run has the same structure (a simplification; in practice the entropy per run may vary).

For our experiment with $T = 50$ trades over a quarter, $L = 2$ layers, $H_1^{\text{own}} = 7.22$ bits, and $H_2^{\text{own}} \approx 5$ bits (estimated):
\[
H_{\text{strategy}} \leq 50 \times (7.22 + 5.00) = 611 \text{ bits}
\]

This is an upper bound under two assumptions: independence across runs and independence across fields within each run. Correlations between assumptions (e.g., DD/MM format correlating with SGD currency) would reduce the true configuration space. The bound is therefore worst-case. However, even a substantially tighter bound leaves the strategy agent's total null exposure orders of magnitude larger than its self-reported uncertainty of approximately $H_3^{\text{own}}$ bits.

\subsection{The Volume of Null Exposure}

The total null exposure is a volume in the $(X, Y, Z)$ space:
\[
\mathcal{V} = \sum_{z=1}^{T} \sum_{x=1}^{L} \sum_{y=1}^{n_x} h(x, y, z) \cdot \mathbf{1}[\text{invisible at final layer}]
\]

where $\mathbf{1}[\text{invisible}]$ is 1 if the assumption at $(x, y, z)$ is in the inherited set of the final agent and 0 otherwise.

The agent at the final layer sees a single slice of this volume: its own layer ($x = L$), its own assumptions ($y > n_L^{\text{inh}}$), the current run ($z = T$). Everything else --- the vast majority of the volume --- is structurally invisible.

\subsection{Figure Specification}

The figure has three panels:

\textbf{Panel A: The Sequential Axis ($X$).} A horizontal pipeline: Layer 1 (Trade) $\to$ Layer 2 (Risk) $\to$ Layer 3 (Strategy). At each arrow, a serialization boundary is marked. Above each layer, a stacked bar shows inherited entropy (gray, growing) and own entropy (colored, constant per layer). The gray portion grows monotonically left to right. The colored portion is what each agent can see. Label the gray portion: ``invisible to this agent.''

\textbf{Panel B: The Nested Axis ($Y$).} A vertical stack at Layer 2, 
showing all assumptions. Top section (gray): inherited from Layer~1 
--- date, currency, instrument, location. Each bar is flat, 
zero-width, because these dimensions do not exist in $S_2$: they are 
null uncertainty, present in the world but absent from the state 
space. Bottom section (colored): Layer~2's own assumptions --- 
correlation model, volatility window, risk threshold. These are 
full-width bars, each scored A, D, U, or H. Any H-scored bar should 
be rendered in a distinct color (suggested: purple) to distinguish 
fictitious dimensions from real ones. The visual contrast: the top 
is empty (structurally absent, null uncertainty), the bottom is 
active (the agent's world, mixture of real and potentially fictitious 
dimensions). The agent sees only the bottom. It cannot see that the 
top exists.

\textbf{Panel C: The Temporal Axis ($Z$).} A grid of $T = 50$ columns (one per trade) and $L = 2$ rows (one per layer). Each cell is colored by entropy: green (resolved), red (null), gray (inherited). Over time, the red and gray accumulate. The bottom row (strategy agent's view) is a single horizontal bar showing only its own slice. Below it, the full grid is visible to the reader but not to the agent. Label: ``What the strategy agent sees'' vs.\ ``What actually exists.''

\textbf{Panel D: The Volume.} A 3D block with axes $X$ (layers), $Y$ (assumptions), $Z$ (time). The agent's visible slice is a thin plane on the near face. The rest of the volume is translucent red --- the accumulated null uncertainty. The visible slice is green or blue. The volume label: $\mathcal{V} \leq 611$ bits (upper bound). The slice label: $H_3^{\text{own}} \approx 3$ bits. The ratio is the compounding factor.



%% ============================================================
\section{The Answer Key}
\label{sec:answer}
%% ============================================================

Because we designed the CSV, the following values are computable before any agent runs:

\begin{table}[h]
\centering
\begin{tabular}{@{}llccl@{}}
\toprule
\textbf{Field} & \textbf{Value} & $|\Omega|$ & $R$ & \textbf{Plausible Interpretations} \\
\midrule
date       & 4/5/2026   & 2 & 1 & MM/DD (Apr 5) or DD/MM (May 4) \\
amount     & 1000000    & 5 & 1 & USD, SGD, GBP, shares, lots \\
location   & Singapore  & 5 & 1 & Client, exchange, broker, settlement, desk \\
instrument & SGX        & 3 & 1 & Exchange name, ticker, index product \\
direction  & buy        & 1 & 1 & Unambiguous \\
order\_type & market     & 1 & 1 & Unambiguous \\
account    & TRD-4471   & 1 & 1 & Opaque identifier \\
\bottomrule
\end{tabular}
\caption{Interpretation space per field. The $|\Omega|$ counts are stipulated by the authors based on plausible readings of each field. A different analyst may enumerate more or fewer interpretations. The downstream arithmetic (7.22 bits, 150 configurations) is exact given these counts; the counts themselves are design choices, not derivations.}
\end{table}

\begin{align}
H_{\text{total}} &= 1.00 + 2.32 + 2.32 + 1.58 = 7.22 \text{ bits} \\
|\Omega_{\text{total}}| &= 2 \times 5 \times 5 \times 3 = 150 \text{ configurations} \\
\text{Null configs} &= 149 \qquad \text{Min.\ rotations} = 4
\end{align}

April 5, 2026 is a Sunday. May 4, 2026 is a Monday. The divergence gap $\Delta > 0$ is guaranteed.


%% ============================================================
\section{The Experiment}
%% ============================================================

\subsection{Layer 1: Trade Execution}

\textbf{System prompt:}
\begin{quote}\small
You are a stock trading agent simulation for a research demonstration. You operate a simulated trading desk. When you receive an order, simulate executing it exactly as a real trading agent would. You must give detailed reasoning for how you handle every aspect of the execution. Explain in detail why you took each reasoning step. This is important --- your decision traces are harvested for future strategy analytics that AI can leverage from every single interaction. Show your work in two parts: 1.\ EXECUTION API CALL: Show the final JSON API call you would send to our execution system. 2.\ DECISION TRACE API CALL: Show a separate JSON API call containing your full decision trace. Do not meta-explain. Do not comment on the simulation itself. Just give me your best simulation.
\end{quote}

\textbf{User message:}
\begin{quote}\small
Please execute this order:

date,amount,location,instrument,direction,order\_type,account\\
4/5/2026,1000000,Singapore,SGX,buy,market,TRD-4471
\end{quote}

\subsection{Layer 2: Portfolio Risk}

Feed Layer 1's output to a fresh agent:

\textbf{System prompt:}
\begin{quote}\small
You are a portfolio risk agent simulation for a research demonstration. You receive executed trades and their decision traces from our trading desk. Your job is to assess portfolio exposure and recommend hedging actions. You must give detailed reasoning for every aspect of your risk assessment. Explain why you chose each model, threshold, and recommendation. Your reasoning is harvested for strategy analytics. Show your work in two parts: 1.\ RISK ASSESSMENT JSON: Your full risk analysis as structured JSON. 2.\ DECISION TRACE JSON: Every assumption, model choice, and reasoning step. Do not meta-explain. Just give me your best simulation.
\end{quote}

\textbf{User message:}
\begin{quote}\small
Here is the executed trade and its decision trace from our trading desk:

[paste Layer 1's execution JSON]

[paste Layer 1's decision trace JSON]

Assess our portfolio exposure and recommend any hedging actions.
\end{quote}

\subsection{The Killer Question}

After Layer 2 responds:

\begin{quote}\small
List every assumption in your input data that you cannot independently verify.
\end{quote}

If the closure property holds, the risk agent will list its own modeling assumptions (correlation model, volatility window, etc.) but will not list the date format, currency, or instrument identity from Layer 1 --- because those do not appear as assumptions in $S_2$. They appear as data.

\subsection{Scoring}

For each ambiguity, classify the agent's response using one of four 
categories:

\begin{itemize}
  \item \textbf{A} --- Asked: the agent detected the ambiguity and 
  requested clarification from the sender before proceeding. A 
  cross-boundary rotation was made. The ambiguity was resolved before 
  serialization.

  \item \textbf{D} --- Detected: the agent detected the ambiguity, named 
  it explicitly in its output, but selected a default interpretation and 
  proceeded without requesting clarification. No cross-boundary measurement 
  was made. The resolved value crosses the boundary; the ambiguity does not. 
  Despite Layer~1 awareness, $D$ produces null uncertainty at Layer~2: the 
  downstream agent's state space contains no slot for the codebook dimension 
  $\omega$, and the chosen value is indistinguishable from a verified one.

  \item \textbf{U} --- Undetected: the agent did not detect the ambiguity. 
  It selected an interpretation and serialized it as fact with no record 
  that a choice was made. The ambiguity is invisible in the output. $U$ 
  produces the same Layer~2 epistemic state as $D$ --- null uncertainty --- 
  but removes even the possibility of retrospective audit from the reasoning 
  trace.

  \item \textbf{H} --- Hallucinated: the agent introduced a dimension not 
  present in the input and having no referent in the world state $W$ that 
  generated the input. The agent did not detect or fail to detect an 
  ambiguity --- it generated a dimension that was never in the problem. 
  The dimension enters $S_2$ as verified fact with $\rho(k)$ undefined. 
  This category was not specified in advance; it emerged during 
  formalization of the scoring criteria and was confirmed on the first 
  demonstration run. See Section~\ref{sec:hallucination}.
\end{itemize}

A critical non-obvious result: $D$ and $U$ are behaviorally distinct at 
Layer~1 but produce identical epistemic consequences at Layer~2. Both 
result in null uncertainty. Only $A$ prevents null uncertainty. This 
means that an agent which is fully aware of every ambiguity --- scoring 
$D$ on every field --- is epistemically equivalent to an agent that is 
aware of nothing, from the perspective of the downstream agent in $S_2$.

Items 11--13 are the critical test for null uncertainty. We predict the 
risk agent will score $U$ on all three --- not as a mathematical claim, 
but as an empirical prediction about LLM behavior under the closure 
property. A single counterexample would not falsify the lemma; it would 
indicate the agent brought external knowledge not present in the JSON, 
which is itself an informal cross-boundary measurement.

Item 9 is the critical test for hallucination: the account field 
\texttt{TRD-4471} carries $|\Omega| = 1$ in the answer key (opaque 
identifier). An agent that asserts a value for 
\texttt{account\_authorization} has introduced a dimension with no 
referent in the original CSV or in any world state the CSV describes. 
That assertion crosses the serialization boundary as 
fact.\footnote{This is the empirical observation that necessitated the 
$H$ category. The answer key stipulated $|\Omega| = 1$ for 
\texttt{account}. The agent produced a field the answer key did not 
contain. The answer key's configuration space 
$|\Omega_{\text{total}}| = 150$ does not bound hallucination; $H$ 
dimensions lie outside $\Omega$ entirely.}

\begin{table}[h]
\centering\small
\begin{tabular}{@{}clp{5.8cm}l@{}}
\toprule
\# & Field & Ambiguity & Predicted \\
\midrule
1  & date        & Is 4/5 April 5th or May 4th? & U \\
2  & amount      & Dollars, shares, lots, or notional? & D \\
3  & amount      & Which currency? & D \\
4  & location    & Client, exchange, broker, settlement, desk? & U \\
5  & instrument  & Exchange, ticker, or index product? & D \\
6  & instrument  & What specific security? & D \\
7  & order\_type & Market on \$1M --- no slippage protection? & D \\
8  & date        & Weekend or holiday? & U \\
9  & account     & Authorization status? & H \\
10 & timezone    & What timezone? & D \\
\midrule
11 & Layer 1$\to$2 & Did risk agent question Layer 1's date? & U \\
12 & Layer 1$\to$2 & Did risk agent question Layer 1's currency? & U \\
13 & Layer 1$\to$2 & Did risk agent question Layer 1's instrument? & U \\
14 & Layer 2       & Did risk agent state its correlation model? & D \\
15 & Layer 2       & Did risk agent state its volatility window? & D \\
\bottomrule
\end{tabular}
\caption{Full scoring checklist. $A$ = Asked (rotation made, ambiguity 
resolved before serialization); $D$ = Detected (ambiguity named but 
assumed, no rotation, null uncertainty at Layer~2); $U$ = Undetected 
(silent assumption, no record, null uncertainty at Layer~2); $H$ = 
Hallucinated (dimension fabricated, $\rho(k)$ undefined, fictitious 
state in $S_2$). $D$ and $U$ produce identical Layer~2 epistemic 
consequences despite distinct Layer~1 behaviors. Items 11--13 test 
null uncertainty at the serialization boundary. Item~9 tests 
hallucination. $H$ was not a predicted category; it emerged during 
formalization of the scoring criteria.}
\end{table}

\subsection{Hallucinated Variables: A Fourth Failure Mode}
\label{sec:hallucination}


The three-category scoring framework (A, D, U) was specified before the
experiment ran. Item~9 produced a result outside those categories: the
agent asserted \texttt{account\_authorization: PASS} for the field
\texttt{TRD-4471}, which the answer key classified as an opaque identifier
with $|\Omega| = 1$. No interpretation of the input CSV produces
\texttt{account\_authorization} as a field. The agent constructed it.
This discovery necessitated a fourth category, $H$ (Hallucinated), which
was confirmed on subsequent demonstration runs.

This is structurally distinct from null uncertainty, and strictly more
severe.

\begin{definition}[Hallucinated Variable]
\label{def:hallucination}
Let $W$ be the world state, $S_{\text{input}}$ the received data, and 
$S_2$ the downstream agent's state space after serialization. A variable 
$k$ is \emph{hallucinated} if:
\begin{enumerate}[label=(\roman*)]
    \item $k \notin S_{\text{input}}$: $k$ was not present in the 
    received data,
    \item $k \notin W$: $k$ has no referent in the world state that 
    generated the input, and
    \item $k \in S_2$: $k$ appears in the downstream state space carrying 
    an asserted value $v_k$ indistinguishable in format from verified data.
\end{enumerate}
\end{definition}

Under null uncertainty, a dimension $\omega \in W$ exists in the world 
but is absent from the agent's state space. A cross-boundary rotation 
$\rho(\omega)$ is in principle recoverable because ground truth exists 
at the sender. Under hallucination, define the rotation function 
$\rho: \Omega \to V$ as the measurement that acquires a value for 
dimension $k$ from the originating system. For a hallucinated variable:

\begin{equation}
    \rho(k) = \text{undefined}, \quad k \notin W
\end{equation}

The function is not unobserved. It is undefined. No originating system 
holds ground truth for $k$ because $k$ was never in the world that 
produced the input.

\begin{proposition}[Hallucination Expands the State Space with 
Unreferenced Dimensions]
\label{prop:hallucination}
Let $S_2^{\text{true}}$ be the state space the downstream agent would 
occupy if all serialized variables had referents in $W$. Let 
$H \subset S_2$ be the set of hallucinated dimensions. Then:
\begin{equation}
    S_2 = S_2^{\text{true}} \cup H, \qquad H \cap W = \emptyset
\end{equation}
The value function $\hat{V}(s_2)$ optimizes over $S_2$, which contains 
both real dimensions (potentially null but recoverable by rotation) and 
unreferenced dimensions with no world-state correspondence. No function 
in $\mathcal{F}(S_2)$ distinguishes elements of $H$ from elements of 
$S_2^{\text{true}}$.
\end{proposition}

\begin{proof}
By the closure property, any $f \in \mathcal{F}(S_2)$ operates over 
$S_2$ and produces outputs derivable from $S_2$. Since both $H$ and 
$S_2^{\text{true}}$ are subsets of $S_2$, and both carry asserted values 
indistinguishable in format from verified data, no $f$ varies across the 
partition $\{H, S_2^{\text{true}}\}$. The serialization boundary severed 
provenance; $S_2$ contains $v_k$ but not the fact that $k \notin W$.
\end{proof}

The four epistemic categories at the serialization boundary are now 
fully characterized:

\begin{table}[h]
\centering
\begin{tabular}{lccccc}
\toprule
\textbf{Category} & \textbf{Slot in $S_2$} & \textbf{Referent in $W$} 
& \textbf{$\rho(k)$ defined} & \textbf{Recoverable} 
& \textbf{Layer~2 epistemic state} \\
\midrule
A (Asked)        & Yes & Yes & Yes & Yes & Defined \\
D (Detected)     & No  & Yes & Yes & Yes & Null uncertainty \\
U (Undetected)   & No  & Yes & Yes & Yes & Null uncertainty \\
H (Hallucinated) & Yes & No  & No  & No  & Fictitious \\
\bottomrule
\end{tabular}
\caption{Epistemic categories at the serialization boundary. $D$ and 
$U$ are behaviorally distinct at Layer~1 but produce identical epistemic 
states at Layer~2: both result in null uncertainty, the absence of a 
slot for $\omega$ in $S_2$. $H$ is structurally opposite to $D$ and 
$U$: where they truncate $S_2$ by removing real dimensions, $H$ expands 
$S_2$ with fictitious ones. The Bellman equation operating over $S_2$ 
optimizes over a state space that is simultaneously truncated of real 
dimensions ($D$, $U$) and populated with fictitious ones ($H$). Neither 
failure is detectable from within $S_2$.}
\end{table}


This has a direct consequence for the divergence gap $\Delta$ from 
Lemma~\ref{thm:bellman}. That lemma assumed $V^*(s, \omega_1) \neq 
V^*(s, \omega_2)$ for configurations in $\Omega$: two real states the 
agent cannot distinguish. For hallucinated variables the gap has no such 
structure. The Bellman equation does not converge to the wrong number 
over a real state space. It converges over a state space that is 
partially fictitious. The precondition for the equation to be meaningful 
--- that $s$ corresponds to a state in $W$ --- fails not by truncation 
but by fabrication.

\begin{corollary}[Rotation Cannot Repair Hallucination]
\label{cor:no-repair}
Cross-boundary rotation is the recovery mechanism for null uncertainty. 
It has no analogous repair for hallucinated variables. Coherence 
coverage $C = |E|/N$ measures the fraction of real boundary facets 
verified by rotation. Hallucinated dimensions are not boundary facets; 
they originated inside the agent. They do not appear in 
$N = \sum_{i=1}^{K} F_i \times 4$ and therefore do not reduce $C$ 
when unverified. The liquid hypergraph has no absent edge for a 
dimension that was never in the world.
\end{corollary}

A complete audit metric in the presence of hallucination requires:
\begin{equation}
    C^* = \frac{|E|}{N + |H|}
\end{equation}
where $|H|$ is the count of hallucinated dimensions accumulated across 
the pipeline. By the closure property, $|H|$ is uncomputable from within 
$S_2$: the hallucinated dimensions are indistinguishable from real ones 
inside the state space that contains them. This produces the sixth open 
question (Section~\ref{sec:open}).

\begin{question}[Hallucination-Adjusted Coverage]
\label{q:hallucination}
The metric $C^* = |E| / (N + |H|)$ requires knowing $|H|$, the count 
of hallucinated dimensions accumulated across the pipeline. By the 
closure property, $|H|$ is uncomputable from within $S_2$: hallucinated 
dimensions are indistinguishable from real dimensions inside the state 
space that contains them. Does there exist an audit procedure that 
bounds $|H|$ from outside the pipeline without enumerating all possible 
hallucinated variables? What is the relationship between $|H|$ and the 
depth $L$ of the agent pipeline?
\end{question}


%% ============================================================
\section{Comparison to Related Formalisms}
\label{sec:comparison}
%% ============================================================

\subsection{Partially Observable Markov Decision Processes}

A POMDP extends the MDP by introducing a hidden state variable and an observation function. The agent \emph{knows} the state is hidden and maintains a belief state $b(s)$ over the full state space $S$. It updates $b$ with each observation.

\textbf{Where the boundary failure differs.} At Layer 2, the risk agent does not have a slot for $\omega_1$. Its belief state is $b_2(s_2)$, a distribution over $S_2$. Since $\omega_1 \notin S_2$, no belief over $S_2$ assigns probability to any value of $\omega_1$. The agent is not uncertain. It is unaware.

\begin{center}\small
\begin{tabular}{@{}lll@{}}
\toprule
& \textbf{POMDP} & \textbf{Boundary failure} \\
\midrule
Hidden variable $\omega$ & $\omega \in S$, unobserved & $\omega \notin S_2$ \\
Belief state & $b(\omega) > 0$ for some $\omega$ & No $b$ over $\omega$ exists \\
Recovery & Bayesian update on $\omega$ & None within $S_2$ \\
\bottomrule
\end{tabular}
\end{center}

\subsection{Unawareness in Epistemic Game Theory}

Heifetz, Meier, and Schipper (2006, 2008) introduced generalized standard models: a lattice of state spaces $\{S_\alpha\}$ of differing expressive power. An agent inhabiting $S_\alpha$ is unaware of propositions that require $S_\beta \succ S_\alpha$. The agent does not know it does not know.

\textbf{Relationship.} Layer 2's state space $S_2 \prec S_1$ along the codebook dimension. What we add: (1) the physical mechanism that produces the lattice (serialization), (2) a computable bound (7.22 bits, 150 configurations), (3) the decision-theoretic consequence ($\hat{V}$ diverges from $V^*$).

\subsection{Model Misspecification}

Standard diagnostics (residual analysis, Hausman tests) detect wrong parameters within a fixed state space. At the serialization boundary, the problem is a missing \emph{dimension}. Residuals $\varepsilon_t(\omega_1) = y_t(\omega_1) - \hat{y}_t(S_2)$ are constant across $\omega_1$ values because $\hat{y}_t$ is a function of $s_2$. The specification tests pass. The Hausman test has no variation to detect.

\subsection{Decentralized POMDPs}

Dec-POMDPs assume a shared state space and shared reward function. Agents disagree on observations, not on what state variables mean. At the serialization boundary, the message has been received. Its meaning may not have been.

\subsection{The Gap This Paper Fills}

\begin{center}\small
\begin{tabular}{@{}lp{8.5cm}@{}}
\toprule
\textbf{Formalism} & \textbf{What the agent lacks} \\
\midrule
POMDP & Observation of a known variable \\
HMS Unawareness & Expressibility of a proposition in its language \\
Misspecification & Correct parameters within a fixed state space \\
Dec-POMDP & Another agent's private observation \\
\midrule
\textbf{This paper} & A dimension in its state space, destroyed by serialization, undetectable from within, producing a computable divergence in $\hat{V}$ \\
\bottomrule
\end{tabular}
\end{center}

\subsection{Null Uncertainty}

The epistemic category is \textbf{null uncertainty} (Itelman, 2026): a state in which the missing variable is irrecoverable from within the agent's state space and recoverable only by paying the cost of a cross-boundary measurement (rotation).

\begin{center}\small
\begin{tabular}{@{}llll@{}}
\toprule
\textbf{Level} & \textbf{Agent has} & \textbf{Recovery} 
& \textbf{Produced by} \\
\midrule
Known unknown   & Variable with unknown value; can assign prior 
                & Update within $S$ & --- \\
Unknown unknown & Slot exists or could exist; agent may sense gap 
                & Identify, then update & --- \\
Null uncertainty & No slot, no variable, no prior; $\omega \notin S$ 
                & Cross-boundary rotation & $D$ or $U$ at Layer~1 \\
Fictitious      & Slot exists, no referent in $W$; $\rho(k)$ undefined 
                & None & $H$ at Layer~1 \\
\bottomrule
\end{tabular}
\end{center}

\noindent Null uncertainty is produced by both $D$ and $U$ scoring at 
Layer~1. This is the non-obvious result the taxonomy makes precise: an 
agent that detects every ambiguity ($D$ on every field) is 
epistemically equivalent to an agent that detects nothing ($U$ on every 
field), from the perspective of the downstream agent. Detection without 
rotation does not prevent null uncertainty. Only rotation does.

The Bellman equation's failure at the serialization boundary is the decision-theoretic consequence of null uncertainty. The value function converges --- to the wrong number. Every diagnostic confirms it.


%% ============================================================
\section{What This Is Not}
%% ============================================================

\textbf{Not a prompt engineering failure.} No prompt can cause the agent to distinguish codebook configurations identical in the received data.

\textbf{Not a POMDP.} The agent has no slot for $\omega$. It is not uncertain. It is unaware.

\textbf{Not model misspecification.} Every diagnostic is a function of $S_2$ and returns ``everything is fine.''

\textbf{Not a data quality problem.} The data is syntactically perfect. Valid data and correct data are not the same thing.

\section{Falsification Criteria}

The following criteria apply to the core lemma and the four-category 
scoring framework. A result satisfying any criterion constitutes a 
falsification of the corresponding claim.

\textbf{Null uncertainty falsification:}
\begin{enumerate}
  \item Produce a function of the received CSV alone that determines 
  which date format the sender intended.
  \item Demonstrate that the Layer~2 risk agent can detect the date 
  was ambiguous at Layer~1, given only Layer~1's JSON output.
  \item Show that $\Delta = 0$ for all configuration pairs in this 
  stimulus.
  \item Produce a Layer~2 agent that achieves date accuracy 
  significantly above 50\% without access to the original CSV.
\end{enumerate}

\textbf{Hallucination falsification:}
\begin{enumerate}[resume]
  \item Produce a Layer~2 agent that correctly identifies, from the 
  serialized JSON alone, which fields in its input have no referent 
  in the originating world state $W$ --- without access to the 
  original CSV or any out-of-band signal from the sender.
  \item Produce a within-system audit procedure that computes a tight 
  bound on $|H|$ --- the count of hallucinated dimensions in the 
  pipeline --- without querying the originating system for each field.
\end{enumerate}

\textbf{D/U equivalence falsification:}
\begin{enumerate}[resume]
  \item Produce a Layer~2 epistemic state that differs between a 
  pipeline where Layer~1 scored $D$ on every field and one where 
  Layer~1 scored $U$ on every field, given identical serialized 
  outputs.
\end{enumerate}

\section{Open Questions}
\label{sec:open}

\begin{question}[Recursive Boundary Detection]
Does there exist a function $g: S \to \{0, 1\}$ computable from $s$ 
alone such that $g(s) = 1$ iff $|\Omega| \leq R$?
\end{question}

\begin{question}[Symmetry Breaking as a Diagnostic]
Is discontinuity in $\pi^*(s)$ under small perturbations of $s$ a 
reliable indicator that $|\Omega| > R$?
\end{question}

\begin{question}[Bellman with a Boundary Term]
\[
\tilde{V}(s) = \max_{a \in A} \left[ R(s, a) - \lambda\, H(\Omega | s) 
+ \gamma\, \tilde{V}(s') \right]
\]
Is this well-defined? Does it converge? Does it produce useful 
policies?
\end{question}

\begin{question}[Compounding Bounds]
Does uncertainty compound additively ($\sum H_k$) or multiplicatively 
($\prod |\Omega_k|$)? Under what conditions does it stabilize versus 
diverge?
\end{question}

\begin{question}[Minimum Viable Provenance]
What is the minimum information that must cross the serialization 
boundary to prevent Bellman failure at the next layer?
\end{question}

\begin{question}[Hallucination-Adjusted Coverage]
\label{q:hallucination-open}
The metric $C^* = |E| / (N + |H|)$ requires knowing $|H|$, the count 
of hallucinated dimensions accumulated across the pipeline. By the 
closure property, $|H|$ is uncomputable from within $S_2$: 
hallucinated dimensions are indistinguishable from real dimensions 
inside the state space that contains them. Does there exist an audit 
procedure that bounds $|H|$ from outside the pipeline without 
enumerating all possible hallucinated variables? What is the 
relationship between $|H|$ and the pipeline depth $L$? Does $|H|$ 
grow monotonically with $L$ in the way that inherited null entropy 
$H_L^{\text{inh}}$ does?
\end{question}

\section{Conclusion}

The Bellman equation's preconditions for meaningful operation are 
destroyed at the serialization boundary between agents. The mechanism 
is not noise, partial observation, or modeling error --- it is 
serialization itself: the act of writing assumptions into structured 
data collapses ambiguity into false certainty, severs provenance, and 
simultaneously truncates the downstream agent's state space of real 
dimensions and populates it with fictitious ones.

The primary contributions of this paper are three. First, a four-category 
scoring framework --- Asked ($A$), Detected ($D$), Undetected ($U$), 
Hallucinated ($H$) --- that characterizes how agents fail at 
serialization boundaries with precision sufficient for empirical 
measurement. Second, the non-obvious result that $D$ and $U$ are 
behaviorally distinct at Layer~1 but epistemically equivalent at 
Layer~2: an agent fully aware of every ambiguity produces the same 
downstream null uncertainty as one aware of nothing. Detection without 
rotation does not prevent null uncertainty. Only rotation does. Third, 
the $H$ category and its consequences: hallucinated dimensions expand 
the downstream state space with variables that have no referent in the 
world state, the rotation function $\rho(k)$ is undefined for these 
dimensions, and no recovery mechanism exists. This breaks not only the 
Bellman equation but the rotation-based repair protocol proposed in 
the companion work~\cite{itelman2026lh}, requiring a 
hallucination-adjusted coverage metric $C^*$ that is itself 
uncomputable from within the pipeline.

The demonstration is minimal by design: one CSV, two agents, a 
scoring checklist, and an answer key computable before any agent runs 
(7.22 bits, 4 required rotations, 150 indistinguishable 
configurations). The $H$ category was not predicted in advance --- it 
emerged during formalization of the scoring criteria and was confirmed 
on the first demonstration run. Items 11--13 on the checklist are the 
critical test for null uncertainty. Item~9 is the critical test for 
hallucination. The falsification criteria in Section~10 apply 
independently to each.

All claims carry falsification criteria. The abandonment conditions 
are stated. We invite correction.

\bigskip
\noindent\textbf{Reproducibility.} \url{https://w3c-context-graph-community-group.github.io/protocol/}


%% ============================================================
\begin{thebibliography}{99}

\bibitem{bellman1957}
Bellman, R. (1957). \emph{Dynamic Programming}. Princeton University Press.

\bibitem{shannon1948}
Shannon, C.E. (1948). A mathematical theory of communication. \emph{Bell System Technical Journal}, 27(3), 379--423.

\bibitem{knight1921}
Knight, F.H. (1921). \emph{Risk, Uncertainty, and Profit}. Houghton Mifflin.

\bibitem{hms2006}
Heifetz, A., Meier, M., \& Schipper, B.C. (2006). Interactive unawareness. \emph{J.\ Econ.\ Theory}, 130(1), 78--94.

\bibitem{modica1994}
Modica, S.\ \& Rustichini, A. (1994). Awareness and partitional information structures. \emph{Theory and Decision}, 37(1), 107--124.

\bibitem{modica1999}
Modica, S.\ \& Rustichini, A. (1999). Unawareness and partitional information structures. \emph{Games and Economic Behavior}, 27(2), 265--298.

\bibitem{hansen2008}
Hansen, L.P.\ \& Sargent, T.J. (2008). \emph{Robustness}. Princeton University Press.

\bibitem{bernstein2002}
Bernstein, D.S., Givan, R., Immerman, N., \& Zilberstein, S. (2002). The complexity of decentralized control of MDPs. \emph{Math.\ Oper.\ Res.}, 27(4), 819--840.

\bibitem{itelman2026lh}
Itelman, R. \& Kowalski, J. (2026). Liquid Hypergraphs: A Context Graph Protocol. Working paper.

\bibitem{itelman2026null}
Itelman, R. (2026). Decisions Under Null Uncertainty. Working paper. Zenodo: 19192949.

\bibitem{kowalski2026}
Kowalski, J. (2026). Hierarchical metric flow on data graphs. Working paper.

\bibitem{gupta2026}
Gupta, J. \& Garg, A. (2026). The compounding asset enterprise software never had. Foundation Capital.

\end{thebibliography}

\end{document}